\chapter{System Architecture}

Nebula separates candidate generation from electrical evaluation. The optimizer
does not estimate circuit performance itself: it submits a parameter vector and
receives a typed evaluation result. This boundary permits PPO, Random Search,
and CEM to share the same simulator and metric definitions.

\begin{figure}[htbp]
  \centering
  \includegraphics[width=\textwidth]{receiver-signal-path.pdf}
  \caption{Implemented receiver signal path. A differential \SI{5}{Gb/s} NRZ
  stimulus is filtered by a four-port channel and applied to the
  transistor-level SKY130 CTLE. The one-tap DFE operates on sampled CTLE output
  values in software. Its feedback subtracts the previous decision, weighted by
  the optimized tap $b_1$, before the current decision is made. The dashed DFE
  boundary is behavioral and must not be interpreted as a transistor schematic.}
  \label{fig:signal-path}
\end{figure}

As shown in \cref{fig:signal-path}, the electrical and behavioral portions have
different validation strength. The CTLE waveform comes from ngspice. Channel
filtering and DFE processing are numerical models. The recovered decisions are
therefore appropriate for design-space exploration but do not include sampler
noise, comparator offset, metastability, feedback delay, or clock jitter.

\section{Software layers}

\begin{description}
  \item[Simulation layer.] \code{simulator/} renders benches, launches ngspice,
  parses outputs, validates S-parameters, filters waveforms, and calculates
  receiver metrics.
  \item[Optimization layer.] \code{rl/} defines parameter grids, state and
  reward schemas, PPO models, checkpoints, and the environment contract.
  \item[Experiment layer.] \code{experiments/} supplies training, evaluation,
  search, qualification, and web-interface entry points.
  \item[Analysis layer.] \code{analysis/} builds specification tables,
  comparisons, qualification summaries, plots, and design catalogs.
\end{description}

The primary receiver API is \code{simulator.receiver.evaluate\_receiver()}.
It returns success status, the first failed stage, combined metrics, per-stage
records, warnings, provenance, runtime, and cache identity.

\section{Traceability and failure handling}

Every expensive evaluation is associated with the circuit parameters,
conditions, channel and port map, model files, ngspice identity, metric schema,
and implementation fingerprints. Hard failures stop later work; execution
failures remain distinguishable from deterministic constraint failures. This is
important because a missing measurement must never be interpreted as a passing
value.

\section{Dashboard and generated deliverables}

The local web interface is a thin controller over the same command-line pipeline
rather than a second evaluator. It provides structured numeric target fields,
backend and PPO-version selection, checkpoint discovery, channel and port-map
configuration, PVT scope, and optional final HD3/noise and eye generation. It
launches each run in a subprocess so a simulator or numerical failure does not
silently corrupt the server state.

Recorded events can be transformed into a bounded SVG run graph whose nodes
represent configuration, targets, checkpoints, candidates, simulator stages,
cache origins, failures, selection, qualification conditions, and artifacts.
The browser view is intentionally abbreviated for large runs; downloadable JSON
retains the complete record. Per-run plots cover reward, strict success, losses,
entropy, electrical metrics, failure stages, and evaluation latency. Missing
measurements appear as gaps rather than zero-valued points.

For a successfully selected design, the output set can include the final
parameter vector, strict specification verdicts, raw JSON and JSONL records,
run-graph JSON/SVG, a parameterized SPICE netlist, a schematic representation,
and CTLE/behavioral-DFE eye artifacts. These outputs are the appropriate judging
interface because they connect a visible result to machine-readable evidence.
